% =====================================================================
% SearchCop — Multi-Camera Person Search via LLM Agent
% Target venue: CVPR 2027 (primary) / ICCV 2027 / NeurIPS 2026
% Build:
%   pdflatex main.tex
%   bibtex main
%   pdflatex main.tex && pdflatex main.tex
%
% NOTE: To switch to the official CVPR style:
%   1. Drop `cvpr.sty` next to this file (from CVPR template release).
%   2. Replace the \documentclass / \usepackage block below with:
%        \documentclass[10pt,twocolumn,letterpaper]{article}
%        \usepackage[review]{cvpr}     % review version
%      or \usepackage{cvpr}            % camera-ready
% =====================================================================

\documentclass[10pt,twocolumn,letterpaper]{article}

% ---- generic packages (CVPR template includes most of these) ----
\usepackage{times}
\usepackage{epsfig}
\usepackage{graphicx}
\usepackage{amsmath}
\usepackage{amssymb}
\usepackage{booktabs}
\usepackage{multirow}
\usepackage{xcolor}
\usepackage{url}
\usepackage[breaklinks=true,bookmarks=false]{hyperref}
\usepackage{cite}

% Margins (only used until cvpr.sty is dropped in)
\usepackage[a4paper,margin=0.7in]{geometry}

% ---- bookkeeping macros ----
\newcommand{\method}{\textsc{SearchCop}\xspace}
\newcommand{\todo}[1]{\textcolor{red}{[TODO: #1]}}
\newcommand{\note}[1]{\textcolor{blue}{[NOTE: #1]}}
\usepackage{xspace}

% =====================================================================
\title{\method: Multi-Camera Person Search as an LLM Agent Task}

\author{Anonymous CVPR submission\\
Paper ID xxxx}

\begin{document}
\maketitle

\begin{abstract}
Multi-camera person search asks: \emph{find the same person as the query
across a network of cameras}. The dominant pipeline today is a single-shot
feature-matching procedure that fuses gait and re-identification (Re-ID)
features in a fixed manner, regardless of context.
We argue that person search is more naturally framed as a
\emph{tool-using decision process}: which camera should be inspected first?
which encoder is reliable for this clip? when should the system fall back
or expand its search? We propose \method, an LLM-driven agent that calls a
library of vision tools (gait encoder, Re-ID encoder, vision-language
captioner, spatio-temporal filter, FAISS retriever, multi-vector reranker,
and an evidence fuser) to answer multi-camera person-search queries.
\method{} produces an auditable \emph{reasoning chain} alongside its Top-K
predictions. We instantiate \method{} on the MEVID benchmark and on three
clothing-change / video Re-ID datasets, and evaluate it under three new
metrics in addition to standard mAP and CMC: \emph{Reasoning Faithfulness},
\emph{Tool-call Efficiency}, and \emph{Cross-camera Recall}.
\todo{plug in numbers from results/summary/all\_runs.xlsx after the
experiment matrix completes}.
\end{abstract}

% =====================================================================
% =====================================================================
% Introduction (CVPR style, 5 paragraphs, ~1 column)
% Structure:
%   P1. Motivation: why multi-camera person search matters in the wild.
%   P2. Gap: classical pipelines are single-shot, modality-fixed, and
%       opaque; clothing change, occlusion, and cross-camera transit
%       expose the limit of one-shot matching.
%   P3. Insight: person search is a multi-step tool-using decision
%       problem; LLM agents excel at exactly that.
%   P4. Approach: SearchCop = Planner / Executor / Critic over a tool
%       library + spatio-temporal prior + self-correcting loop, with
%       five enumerated contributions.
%   P5. Results headline + roadmap of the paper.
% =====================================================================
\section{Introduction}
\label{sec:intro}

Locating a person of interest across a network of surveillance cameras is a
foundational capability for public safety, missing-person investigation,
and large-scale forensic analysis. In contrast to image-level person
re-identification (Re-ID), the operational task analysts actually face is
\emph{multi-camera person search}: given a short query (an image, a video
clip, or even a textual description), determine \emph{when} and \emph{from
which camera} the same individual was seen, and return a ranked list of
candidate tracklets together with an explanation that a human can verify.
Modern monitoring deployments comprise hundreds of overlapping or sparsely
connected cameras, and queries arrive with heterogeneous quality: a frontal
crop here, a back-view tracklet there, sometimes only a witness sentence.
This setting demands a search system that is not only accurate, but also
\emph{adaptive}, \emph{cost-aware}, and \emph{auditable}.

The dominant approach today reduces multi-camera search to a single-shot
feature-matching procedure: encode the query with a Re-ID
network~\cite{he2021transreid,li2023clipreid} or a gait
network~\cite{chao2019gaitset,fan2023opengait}, retrieve top-$K$ candidates
from a fixed-modality FAISS index, and optionally late-fuse the two
modalities with hand-tuned weights. This recipe ignores three structural
properties of the problem. First, query quality is highly variable -- the
right encoder for a clean walking sequence is the wrong encoder for a
heavily occluded crop, yet existing systems commit to a modality at
indexing time. Second, multi-camera networks have rich spatio-temporal
priors: a person seen at camera~A at time $t$ cannot physically appear at
a non-adjacent camera~B within a few seconds, yet most retrieval pipelines
discard topology entirely~\cite{davila2023mevid,zheng2017prw}. Third,
clothing-change benchmarks~\cite{yang2019prcc,qian2020ltcc} reveal that any
single modality, however carefully fine-tuned, fails on a non-trivial slice
of queries; the fix is not a stronger encoder, but a \emph{decision
policy} that knows when to switch.

Our key insight is that multi-camera person search is more naturally cast
as a \emph{multi-step tool-using decision problem} than as a one-shot
retrieval problem. Recent advances in tool-using large language model
(LLM) agents -- from ReAct~\cite{yao2023react} and
Toolformer~\cite{schick2023toolformer} to Visual
Programming~\cite{gupta2023visprog} and Chameleon~\cite{lu2023chameleon}
-- show that LLMs, when given a structured tool library, can plan
non-trivial sequences of vision/database calls and produce traceable
rationales. We bring this paradigm to person search. Crucially, the same
encoders that prior work uses as monolithic pipelines (gait, Re-ID,
vision-language captioners) become \emph{callable tools} in the LLM's
action space, and the same camera-topology priors that are usually
discarded become a \emph{spatio-temporal filter tool}. The LLM's job is
not to look at pixels: it is to decide \emph{which tool to call next}
given the current evidence, and to declare \emph{when} the answer is
trustworthy.

Building on this insight we propose \method, the first LLM agent for
multi-camera person search. \method{} is built around a Planner /
Executor / Critic loop driven by a frontier LLM. The Planner emits one
function call at each step, choosing among $10$ tools that wrap the gait
encoder, the CLIP-Re-ID encoder, a multi-vector reranker, a vision-language
captioner, a spatio-temporal filter, a quality scorer, a FAISS coarse
retriever, and an evidence fuser; the Executor runs the tool and returns a
typed result; the Critic scores the step's usefulness and the running
confidence, triggering a \emph{self-correcting} replan whenever
modalities disagree or confidence stagnates. We evaluate
\method{} on MEVID~\cite{davila2023mevid}, MARS~\cite{zheng2016mars},
DukeMTMC-VideoReID~\cite{wu2018duke}, and the clothing-change benchmarks
PRCC~\cite{yang2019prcc} and LTCC~\cite{qian2020ltcc}, using both standard
mAP/CMC and three new metrics tailored to agent-based retrieval:
\emph{Reasoning Faithfulness}, \emph{Tool-call Efficiency}, and
\emph{Cross-camera Recall}. To summarize, our contributions are:

\begin{itemize}\setlength\itemsep{2pt}
  \item \textbf{Task reformulation.} We are the first to cast multi-camera
        person search as a \emph{tool-using LLM agent task}, replacing the
        prevailing single-shot matching paradigm with an explicit
        plan--act--criticize loop.
  \item \textbf{Tool library and schema.} We design a $10$-tool library
        (gait, Re-ID, VLM caption, spatio-temporal filter, quality scorer,
        FAISS, multi-vector reranker, evidence fuser, SAM silhouette,
        track extractor) with strict JSON schemas, so the Planner can
        compose them safely without function-calling support from the
        underlying LLM gateway.
  \item \textbf{Spatio-temporal reasoning prompt.} We integrate the
        camera-topology graph and time-window prior into the Planner
        prompt, so the agent prunes physically infeasible candidates
        before invoking heavy encoders.
  \item \textbf{Self-correcting agent loop.} We introduce a Critic that
        outputs a calibrated confidence and a \emph{needs-correction}
        flag, enabling the system to expand cameras, broaden time
        windows, or switch encoders when evidence is weak. We further
        prevent pathological loops with a tool-argument deduplication
        guard.
  \item \textbf{New evaluation protocol.} Beyond the standard mAP/CMC, we
        propose \emph{Reasoning Faithfulness} (alignment between stated
        rationale and actual tool output) and \emph{Tool-call Efficiency}
        (average number of calls per correctly answered query), turning
        the otherwise opaque agent behaviour into a measurable artefact.
\end{itemize}

Empirically, \method{} achieves \todo{X.X}\% mAP on MEVID, an absolute
improvement of \todo{$+\Delta$}\% over the strongest hand-tuned
gait+Re-ID fusion baseline, while also raising rank-1 by \todo{X.X}\% on
PRCC and LTCC. The agent makes on average \todo{N} tool calls per query,
\todo{Y}\% of which are validated by the Critic as useful; the resulting
reasoning chains are rated as faithful by \todo{Z}\% of human evaluators.
The remainder of the paper is organized as follows.
Section~\ref{sec:related} reviews video Re-ID, multi-camera person search,
and tool-using LLM agents. Section~\ref{sec:method} describes the tool
library, the Planner / Critic prompts, and the self-correcting loop.
Section~\ref{sec:exp} reports our experiments and ablations. We conclude
in Section~\ref{sec:conclusion}.


% --- placeholders for next sections (will be filled D10--D13) ---
% =====================================================================
% Related Work (CVPR style, ~1 column, 4 subsections)
% Each subsection ends with one sentence that states the gap and how
% SearchCop addresses it. The final paragraph contrasts SearchCop with
% the closest tool-using vision-language agents.
% =====================================================================
\section{Related Work}
\label{sec:related}

\subsection{Video Re-Identification}
\label{ssec:rel_videoreid}
The dominant body of work on person retrieval treats the problem as
representation learning over single-camera tracklets. Image-level methods
such as TransReID~\cite{he2021transreid} and CLIP-ReID~\cite{li2023clipreid}
introduced strong transformer or vision-language pretraining backbones,
while video-level methods such as AP3D~\cite{gu2020ap3d},
GRL~\cite{liu2021grl}, and STMN~\cite{eom2021stmn} aggregate temporal
information through 3D convolutions, attention, or memory-based pooling.
Standard benchmarks include MARS~\cite{zheng2016mars} and
DukeMTMC-VideoReID~\cite{wu2018duke}; both fix the gallery to a single
modality and assume the system retrieves Top-$K$ in one shot.
\textbf{Gap.} These approaches optimize a single similarity function and
deliver no \emph{decision policy}: the system cannot, for example, decide
to switch modality when the query is occluded, nor explain why a Top-1
prediction was chosen. \method{} preserves these encoders as
\emph{tools} but lifts modality selection and aggregation into an explicit
LLM-driven plan.

\subsection{Multi-Camera Person Search and Spatio-Temporal Reasoning}
\label{ssec:rel_mcps}
Multi-camera person search jointly tackles detection and re-identification
across multiple cameras, with PRW~\cite{zheng2017prw} and
CUHK-SYSU~\cite{xiao2017cuhksysu} as canonical benchmarks; recent
MEVID~\cite{davila2023mevid} explicitly targets multi-environment,
long-time-span video. A handful of works exploit
\emph{spatio-temporal} priors -- camera topology, transit-time histograms,
or graph propagation -- to refine retrieval~\cite{lv2018global,wang2019stp}.
Yet, both lines treat each gallery item independently and never plan a
\emph{sequence} of retrieval actions; topology, when used at all, is
applied as a fixed post-filter rather than a tool the system can choose
to invoke. \textbf{Gap.} The decision of \emph{when} to consult topology,
\emph{which} encoder to apply, and \emph{whether} the answer is good
enough is left to manual heuristics. \method{} promotes each of these
decisions to first-class actions in an LLM agent's tool space.

\subsection{Clothing-Change Re-ID and Gait Recognition}
\label{ssec:rel_cc_gait}
When clothing changes between camera views, appearance-based features
become unreliable. Two complementary lines emerged. (i) Clothing-aware
representation learning, exemplified by CAL~\cite{gu2022cal} on
PRCC~\cite{yang2019prcc} and LTCC~\cite{qian2020ltcc}, learns features
adversarially insensitive to clothes. (ii) Gait recognition exploits
walking dynamics rather than appearance: GaitSet~\cite{chao2019gaitset},
GaitPart~\cite{fan2020gaitpart}, and the OpenGait
family~\cite{fan2023opengait} demonstrate strong cross-view performance
on silhouette inputs. Both lines, however, presume an \emph{a priori}
choice of modality at indexing time, and offline late fusion
typically falls back to fixed weighted sums.
\textbf{Gap.} Real surveillance queries rarely come with a clean modality
hint -- a frontal occluded crop calls for Re-ID, a clean back-view
tracklet calls for gait. \method{} keeps both encoders as on-demand
tools and lets the Planner decide which one to trust per query, with the
Critic invalidating the call when modality disagreement is detected.

\subsection{LLM Agents and Tool Use}
\label{ssec:rel_agents}
Tool-using LLM agents are a fast-growing line of work originating from
ReAct~\cite{yao2023react} and Toolformer~\cite{schick2023toolformer},
which interleave reasoning and acting via natural-language traces and
self-supervised tool annotation, respectively. Building on these,
HuggingGPT~\cite{shen2023hugginggpt} orchestrates Hugging Face models for
multi-modal tasks; Visual Programming~\cite{gupta2023visprog} and
Chameleon~\cite{lu2023chameleon} compose vision tools into programs for
single-image VQA and reasoning. Beyond, AgentBench~\cite{liu2024agentbench}
and ToolBench~\cite{qin2024toolbench} measure agent capability on coding,
web, and API tasks. \textbf{Gap.} Existing agent benchmarks focus on
\emph{capability} (can the agent solve a task?) rather than
\emph{retrieval faithfulness} (does the agent's reasoning match its tool
output?), and none of these systems tackle large-scale multi-camera
gallery retrieval. \method{} adapts the tool-using paradigm to person
search and additionally proposes \emph{Reasoning Faithfulness} and
\emph{Tool-call Efficiency} as new evaluation primitives that future
agent-based retrieval systems can adopt.

\subsection{Vision-Language Agents for Video and Surveillance}
\label{ssec:rel_vlmagents}
Recent video-centric multimodal LLMs such as
VideoChat~\cite{li2023videochat}, Video-LLaMA~\cite{zhang2023videollama},
Video-ChatGPT~\cite{maaz2023videochatgpt}, and
TimeChat~\cite{ren2024timechat} push LLMs to understand long videos via
captioning or temporal QA. Closer to retrieval, AssistGPT~\cite{gao2023assistgpt}
plans modality-specific calls for video QA, and a small literature on
LLM-driven surveillance studies anomaly description. \textbf{Gap.} These
systems address \emph{describing} or \emph{answering questions about} a
single video, not \emph{searching} a multi-camera gallery for the
\emph{same person} under spatial-temporal constraints; they also have no
notion of self-correcting retrieval. To our knowledge, \method{} is the
first LLM agent that targets this exact problem.

\paragraph{Positioning.} Table~\ref{tab:positioning} contrasts \method{}
with the most relevant prior work along five axes: multi-camera support,
gait+Re-ID modality fusion, explicit decision making, interpretable
output, and evaluation on MEVID. \method{} is the first system to be
positive on all five.

\begin{table*}[t]
\centering
\small
\caption{Positioning of \method{} against representative prior work.
\textbullet\ = supported, $\circ$ = partial, blank = unsupported.}
\label{tab:positioning}
\begin{tabular}{lccccc}
\toprule
Method & Multi-cam & Modality fusion & Decision making & Interpretable & MEVID eval \\
\midrule
TransReID~\cite{he2021transreid}              & $\circ$ &        &        &        &        \\
CLIP-ReID~\cite{li2023clipreid}                & $\circ$ & text aux &      & $\circ$ &        \\
CAL~\cite{gu2022cal} (clothing-change)         &         & clothes-aware &    &        &        \\
OpenGait~\cite{fan2023opengait} family         &         & gait only &      &        &        \\
HuggingGPT~\cite{shen2023hugginggpt}           &         & LLM-tool & \textbullet & \textbullet &  \\
Visual Programming~\cite{gupta2023visprog}     &         & LLM-tool & \textbullet & \textbullet &  \\
Chameleon~\cite{lu2023chameleon}               &         & LLM-tool & \textbullet & \textbullet &  \\
AssistGPT~\cite{gao2023assistgpt}              & $\circ$ & LLM-tool & \textbullet & \textbullet &  \\
\midrule
\textbf{\method{} (ours)} & \textbullet & gait+Re-ID+VLM & \textbullet & \textbullet & \textbullet \\
\bottomrule
\end{tabular}
\end{table*}

% Day 11: Method
% Source material: docs/architecture.md
\section{Method}
\label{sec:method}

\todo{draft Day 11 — overview figure; tool library + JSON schema; Planner/Executor/Critic prompts; self-correcting loop; spatio-temporal reasoning; Algorithm 1 (SearchCopLoop); Algorithm 2 (EvidenceFuse).}

% Day 12: Experiments
% Source material: results/summary/all_runs.xlsx
\section{Experiments}
\label{sec:exp}

\todo{draft Day 12 — datasets (MEVID, MARS, DukeMTMC-VideoReID, PRCC, LTCC); metrics (mAP, CMC, Reasoning Faithfulness, Tool-call Efficiency, Cross-camera Recall); main table; ablations (tools, CoT, self-correct, LLM tier); qualitative cases.}

% Day 13: Conclusion
\section{Conclusion}
\label{sec:conclusion}

\todo{draft Day 13 — recap contributions, limitations (latency, single-LLM dependency), future work (extending to face / vehicle / scene attribute search; on-device deployment).}


% =====================================================================
\bibliographystyle{ieee_fullname}
\bibliography{references}

\end{document}
